\documentclass[12pt,a4paper]{report}

% =============================================================================
%  Lucrare de disertatie -- Cliq
%  Fisier unic (conform alegerii): main.tex + capitole ar putea fi separate
%  ulterior prin \input{...} daca lucrarea creste mult in dimensiune.
% =============================================================================

\usepackage[utf8]{inputenc}
\usepackage[T1]{fontenc}
\usepackage[romanian]{babel}
\usepackage{lmodern}
\usepackage{microtype}
\usepackage[a4paper,left=3cm,right=2.5cm,top=2.5cm,bottom=2.5cm]{geometry}
\usepackage{setspace}
\usepackage{graphicx}
\usepackage{amsmath,amssymb}
\usepackage{enumitem}
\usepackage{csquotes}
\usepackage[colorlinks=true,linkcolor=black,citecolor=black,urlcolor=blue,bookmarksnumbered=true]{hyperref}

\onehalfspacing

\newcommand{\titlullucrarii}{Cliq, Generarea automată a claselor Page Object pentru testarea automată a interfețelor web}

\begin{document}

% =============================================================================
%  Pagina de titlu
% =============================================================================
\begin{titlepage}
    \centering
    \vfill
    {\Huge\bfseries \titlullucrarii\par}
    \vfill
\end{titlepage}

\pagenumbering{roman}
\tableofcontents
\newpage
\pagenumbering{arabic}

% =============================================================================
\chapter{Introducere}
% =============================================================================

\section{Context și motivație}

Testarea automată a interfețelor web se bazează, în cea mai mare parte, pe
Selenium WebDriver și pe patternul \emph{Page Object Model} (POM), care
izolează localizarea elementelor și interacțiunile cu acestea în clase
dedicate, separat de logica scenariilor de test. Scrierea și mentenanța
acestor clase rămân însă activități manuale: fiecare pagină nouă sau
modificată impune identificarea elementelor din DOM, alegerea unor locatori
stabili și implementarea metodelor de interacțiune, iar orice discrepanță
între cod și structura reală a paginii produce teste \emph{flaky}, care
eșuează intermitent din cauza unui locator devenit invalid, nu a unui defect
real al aplicației.

Cliq, aplicația care face obiectul prezentei lucrări, pornește de la
observația că această informație, elementele interactive ale unei pagini,
atributele lor stabile, tipul de interacțiune așteptat, poate fi extrasă
automat prin scanarea paginii reale, în loc să fie transcrisă manual din
inspectorul browserului.

\section{Problema abordată}

Instrumentele existente de generare a codului de test au, în general, patru
limitări: regenerarea suprascrie orice personalizare adăugată manual în clasa
generată; rescanarea unei pagini deja cunoscute aplică automat diferențele
detectate (elemente noi sau dispărute), fără control explicit al
dezvoltatorului; codul generat urmează convenții fixe, diferite de cele deja
adoptate în proiectul în care este integrat; iar scanarea, generarea și
rularea testelor rămân activități separate de mediul de dezvoltare în care
codul este de fapt scris și revizuit.

\section{Obiectivele lucrării}

Pornind de la aceste limitări, lucrarea își propune proiectarea și
implementarea unei aplicații care automatizează generarea claselor Page
Object, tratând regenerarea codului ca pe un proces sigur și transparent
pentru dezvoltator. Obiectivele concrete sunt:

\begin{enumerate}[label=O\arabic*., noitemsep, topsep=4pt]
    \item scanarea automată a unei pagini web și atribuirea de locatori
    stabili elementelor identificate, în ordine de prioritate explicită;
    \item generarea claselor Page Object în C\# și a testelor NUnit
    corespunzătoare, executabile și independent, prin \texttt{dotnet test};
    \item un mecanism de regenerare care păstrează metodele scrise manual și
    cere o decizie explicită pentru fiecare, în locul suprascrierii tăcute;
    \item un flux de tip \emph{staged review} pentru diferențele detectate la
    rescanarea unei pagini, supuse acceptării sau respingerii explicite;
    \item preluarea convențiilor de denumire dintr-un page object existent,
    oferit ca exemplu, astfel încât codul generat să se integreze stilistic
    în proiect;
    \item integrarea instrumentului într-o extensie pentru Visual Studio, pe
    lângă aplicația desktop de sine stătătoare construită pe Electron.
\end{enumerate}

\section{Contribuțiile lucrării}

Rezultatul concret este aplicația Cliq: un motor de scanare bazat pe
Puppeteer, capabil să reconcilieze rezultatul cu o scanare anterioară a
aceleiași pagini; un generator de cod C\#/NUnit, care produce inclusiv o
clasă de bază comună (\texttt{BasePage}) stilizată după un exemplu existent;
un mecanism de fuziune a metodelor personalizate bazat pe atributul
\texttt{[CliqCustomMethod]}, care nu permite pierderea tăcută a unei metode;
și două găzduiri, o aplicație desktop Electron și o extensie Visual Studio,
care expun același set de operații prin intermediul unui proces auxiliar
(\emph{sidecar}) autentificat.

% =============================================================================
%  Capitolele urmatoare - schelet, de completat pe masura ce avansam.
% =============================================================================

% \chapter{Stadiul actual al domeniului}
% \chapter{Arhitectura sistemului Cliq}
% \chapter{Implementare}
% \chapter{Testare și evaluare}
% \chapter{Concluzii și direcții viitoare}

% \bibliographystyle{plain}
% \bibliography{bibliografie}

\end{document}
